\documentclass[12pt]{article}
\usepackage[utf8]{inputenc}
\usepackage[T1]{fontenc}
\usepackage{hyperref}
\title{Tayart Tamajeq Language Resources for K-12 Curriculum Development}
\author{Compiled by Claude Sonnet 5 (claude-sonnet-5) for Andr\'as Kornai}
\date{}
\begin{document}
\maketitle
\section{Introduction}
Tayart Tamajeq (thz; Glottolog taya1257) is spoken in Niger by roughly 250,000 people. Omniglot reports that it is written mainly in Latin script, with Tifinagh and Arabic-based scripts also in use, and taught in some primary schools. I found that only on Omniglot and did not verify it elsewhere.
\section{Primary Educational Resources}
No textbooks or primers turned up in the search results.
\section{Linguistic and Reference Materials}
\begin{itemize}
\item Towards an Electronic Dictionary of Tamajaq Language in Niger (Enguehard and Modi, AfLaT 2009): \url{https://aflat.org/node/316}
\item Slides for the same paper, with about 6,378 base entries: \url{https://www.aflat.org/files/aflat2009/slides/tamajaq.pdf}
\item Paper PDF: \url{https://preview.aclanthology.org/moar-dois/W09-0712.pdf}
\item Omniglot page on the writing systems: \url{https://www.omniglot.com/writing/tayarttamajeq.htm}
\end{itemize}
\section{Religious and Community Resources}
None found beyond the audio page below.
\section{Digital and Audio Resources}
Global Recordings Network audio: \url{https://www.globalrecordings.net/en/language/1793}
\section{Curriculum Development Context}
The electronic dictionary project involved Niger's Ministry of National Education. The search results did not show the current state of classroom materials.
\end{document}